\documentclass[11pt,a4paper]{article}

% ---------- Packages ----------
\usepackage[margin=1.15cm]{geometry}
\usepackage{titlesec}
\usepackage{enumitem}
\usepackage{hyperref}
\usepackage{fontspec}

% ---------- Fonts ----------
\setmainfont{Times New Roman}

% ---------- Section Style ----------
\titleformat{\section}
  {\large\bfseries}
  {}
  {0em}
  {}
  [\titlerule]

\titlespacing*{\section}{0pt}{7pt}{5pt}

% ---------- List Style ----------
\setlist[itemize]{noitemsep, topsep=2pt, leftmargin=1.2em}

% ---------- Hyperlink ----------
\hypersetup{
  colorlinks=true,
  urlcolor=black
}

% ---------- Document ----------
\begin{document}
\pagenumbering{gobble}

% ---------- Header ----------
\noindent
\begin{minipage}[t]{0.62\textwidth}
  {\fontsize{22}{24}\selectfont \textbf{Rui Shao}}
\end{minipage}
\hfill
\begin{minipage}[t]{0.36\textwidth}
  \raggedleft
  {\small +86-15623713170 / +81-07090122378}\\
  {\small \href{mailto:sr1054461216@gmail.com}{sr1054461216@gmail.com}}\\
  {\small \href{https://github.com/ryan42xyz}{github.com/ryan42xyz}}
\end{minipage}

\vspace{-2pt}

% ---------- Experience ----------
\section*{Experience}

% ===== Datavisor =====
\noindent
\href{https://www.datavisor.com}{\textbf{Datavisor}} \hfill \textit{Senior Site Reliability Engineer (Japan)} \hfill 2025/10 -- Present \\
\phantom{\href{https://www.datavisor.com}{\textbf{Datavisor}}} \hfill \textit{Site Reliability Engineer (China)} \hfill 2025/03 -- 2025/10

\vspace{4pt}

\noindent\textbf{\textit{Deployment \& Reliability}} \hfill \textit{(\textbf{AWS} \textbf{6 regions} | \textbf{Kubernetes} \textbf{25+ clusters} | \textbf{600+ nodes} )} \hfill \mbox{}
\begin{itemize}
  \item Led a multi-version upgrade of self-managed \textbf{Kubernetes} clusters on \textbf{AWS}, coordinating \textbf{control-plane} upgrades, system component compatibility (\textbf{Calico}, \textbf{AWS cloud-controller-manager}, \textbf{cluster-autoscaler}), and worker-node replacement via \textbf{ASG Launch Template}.
  \item Built a staged upgrade workflow with preflight compatibility checks (APIs, CRDs), etcd snapshot rollback, and controlled node draining with \textbf{PodDisruptionBudgets} to protect stateful workloads (MySQL, Kafka, ClickHouse).
  \item Implemented traffic-safe rollout with blue-green validation and progressive traffic warm-up using \textbf{SLO} (latency, error rate, critical APIsuccess, restarts).
\end{itemize}

\vspace{2pt}

\noindent\textbf{\textit{Traffic \& Multi-Region}}
\begin{itemize}
  \item Designed \textbf{API Gateway}-based cross-cluster traffic control for stateless \textbf{Kubernetes} services, enabling weighted routing between primary and standby clusters for canary releases, cutovers, and incident failover.
  \item Implemented \textbf{SLI}-driven failover logic using service-level signals, with staged failback and warm-up validation to prevent route flapping and ensure stable recovery under production traffic.
\end{itemize}

\vspace{2pt}

\noindent\textbf{\textit{Observability}}
\begin{itemize}
  \item Shifted monitoring from resource-centric alerts to service-level \textbf{SLO} signals (latency, error rate, burn-rate), improving alert quality and reducing noisy infrastructure alarms in a multi-cluster \textbf{Kubernetes} environment.
  \item Built an AI-assisted incident investigation workflow that correlates metrics, logs, deployment history, and service dependencies into a single operator view, accelerating on-call triage and reducing time-to-diagnosis.
\end{itemize}

\vspace{2pt}

\noindent\textbf{\textit{Cost Optimization}}
\begin{itemize}
  \item Reduced cloud spend by optimizing compute models (\textbf{Reserved} baseline, On-Demand scaling, \textbf{Spot} for tolerant workloads) while maintaining reliability guardrails and fallback capacity for production services.
  \item Improved cost efficiency across compute, storage, and network layers by implementing storage lifecycle policies, optimizing log retention, and reducing unnecessary cross-AZ data transfer in a multi-cluster Kubernetes environment.
\end{itemize}

\vspace{2pt}

\noindent\textbf{\textit{On-call Incident Response}}
\begin{itemize}
  \item Led production on-call incident response across Kubernetes and AWS, triaging and diagnosing cross-layer connectivity and performance issues across traffic, networking, application, and node layers using symptom-driven analysis.
  \item Restored service through controlled mitigations such as traffic splitting, selective degradation, bounded scaling, and node replacement while limiting blast radius.
\end{itemize}

\vspace{5pt}

% ===== Intel =====
\noindent
\href{https://www.intel.com}{\textbf{Intel}} \hfill \textit{Cloud Software Development Engineer} \hfill 2022/06 -- 2025/02 (China)

\vspace{2pt}
\begin{itemize}
  \item Built \textbf{Kubernetes} platform tooling using operator-style controllers, IaC (GitOps/Ansible) cluster management.
  \item Developed Go-based \textbf{gRPC} service using \textbf{OpenTelemetry SDK} to transform structured messages into metrics and stream high-throughput time-series via \textbf{Kafka} into \textbf{Prometheus} and \textbf{InfluxDB}.
\end{itemize}

\vspace{4pt}

% ===== Internships =====
\noindent
\href{https://www.tencent.com}{\textbf{Tencent Inc.}} \hfill \textit{Software Development Engineer Intern} \hfill 2021/05 -- 2021/09 \\
Built Go-based gRPC microservices for live-streaming workloads, working with Kafka, MySQL, and Redis.

\vspace{2pt}
\noindent
\href{https://www.bytedance.com}{\textbf{ByteDance Inc.}} \hfill \textit{Software Development Engineer Intern} \hfill 2021/03 -- 2021/05

\vspace{1pt}
\noindent
\href{https://www.data.ai}{\textbf{App Annie}} \hfill \textit{Web Backend Engineer Intern} \hfill 2020/11 -- 2021/02

\vspace{0pt}
\noindent
\href{https://www.baidu.com}{\textbf{Baidu Inc.}} \hfill \textit{Data Engineer Intern} \hfill 2020/01 -- 2020/10

% ---------- Skills ----------
\section*{Skills}
\textbf{Languages:} Python, Golang, Java, Shell, SQL \\
\textbf{Infrastructure \& Data:} AWS, GCP, Kubernetes, Docker, MySQL, Yugabyte, ClickHouse, Redis, Kafka, Elasticsearch \\
\textbf{Observability \& Platform:} Prometheus, InfluxDB, Grafana, Loki, gRPC, Helm, Git, Jenkins

% ---------- Education ----------
\section*{Education}
\textbf{Master}, {\small University of Science \& Technology Beijing, Computer Science and Technology (Beijing, China)} \hfill {\small 2019/09 -- 2022/06}

\vspace{4pt}
\noindent
\textbf{Bachelor}, {\small Wenhua College, Electronic Information Engineering (Wuhan, China)} \hfill {\small 2014/09 -- 2018/06}

\end{document}
